\chapter{\ph{Título do anexo}}
\label{anx:a}

\orientacao{Anexo é o lugar do material que interrompe a leitura mas precisa ficar registrado: tabelas longas de parâmetros e de resultados brutos, lista de componentes com especificação, endereços de repositórios com a versão entregue (\emph{commit} ou \emph{tag}) e registros de propriedade intelectual. Não cole trecho extenso de código; aponte o repositório. Cada anexo é citado no corpo do relatório. Para tabela que atravessa páginas, use \texttt{longtable}, que repete o cabeçalho.}

\ph{Conteúdo do anexo, citado no corpo do relatório.}
